% =============================================================================
% Section 5 results fragment: the Energy Security pillar, final run.
% Every number below is written to paper/results/es.json by scripts/paper/fig_es.R,
% which also draws the five figures and paper/tables/es_leaders.tex.
% Sample: the eleven headline technologies of energy_security_index in index_outputs.rds
% (final rerun of 2026-10-12). Rankings, distributions, fixed-effect decompositions,
% correlations, contribution averages and coalition averages use the rankable cells only
% (sovereign economy, imputed weight share at most 0.5; config/rankability.yml).
% =============================================================================

\paragraph{Coverage and rankability.} The pillar is scored for 6{,}232
country-technology-stage cells of the eleven headline technologies. They cover 219 economies,
195 sovereign states and 24 territories, and 30 of the 33 technology-stage pairs, ten at each
stage; Geothermal upstream and midstream and Electric Grid downstream have no Energy Security
category. No pair is constant across economies or unscored, and no cell rests entirely on
imputed components. The mean imputed weight share of equation~\eqref{eq:imputed_share} is 0.41
upstream, 0.23 midstream and 0.27 downstream, and 38 per cent of the category scores contain at
least one imputed component. Every ranking and summary below uses the rankable cells only: a
cell is rankable if its economy is sovereign and at most $\tau=0.5$ of its index weight is
imputed. The rule keeps 4{,}163 of the 6{,}232 cells. It removes the 607 cells of the 24
territories and 1{,}462 cells of sovereign economies whose score rests mainly on imputation, and
it bites hardest upstream, where 1{,}124 of 2{,}181 cells are rankable, against 1{,}397 of 2{,}002
midstream and 1{,}642 of 2{,}049 downstream. The reason is the grid rule of
Section~\ref{subsec:norm_missing}: the world-level price series places every economy of the
World Bank list in the upstream grid of each technology with a price series, while the
country-specific mineral and fossil-fuel sources cover far fewer economies. Among the rankable
cells the mean imputed weight share falls to 0.18 upstream, 0.08 midstream and 0.22 downstream.
A pair holds between 68 rankable economies (Nuclear downstream) and 195 (Nuclear upstream).
Within each pair every scored economy observes the same set of categories, so a cross-country
comparison within a pair compares one construct; across pairs the bundle ranges from four to
seven categories upstream, one to three midstream and one to four downstream
(Table~\ref{tab:es_structure}). The upstream trade baskets of the final data release give Trade
categories to Electric Vehicles (battery-mineral ores and magnet rare earths), Wind (rare
earths), Solar (silicon metal) and Nuclear (uranium ores and concentrates) upstream, so that nine
of the ten upstream pairs carry a Trade category; Green Hydrogen upstream is the exception.

\paragraph{Distribution across cells.} Figure~\ref{fig:es_distribution} shows the pillar by
technology within each stage. The stages differ in level and in shape. The median rankable cell
scores 0.28 upstream, 0.46 midstream and 0.44 downstream, with interquartile ranges of 0.19 to
0.37, 0.36 to 0.53 and 0.35 to 0.52, and the upstream distribution is strongly right-skewed
(skewness 1.24, against 0.40 midstream and 0.09 downstream). In most upstream pairs the bulk of
economies sits in a narrow band near the bottom while a few lie far above it: in Electric Grid,
Solar and Green Hydrogen upstream the interquartile range is narrower than 0.05. This is an
endowment pattern, filtered through the construction. Upstream, Reserves, Production, Minerals
Trade and Investment are registers of holders, producers, exporters and investors that carry
the zero rule and the zero-anchored S-curve of equation~\eqref{eq:scurve_zero}; in the median
upstream pair, 71 per cent of the rankable economies score zero on Reserves and 84 per cent on
Production. The few economies with mineral or fossil resources therefore separate sharply from
the rest, and the pair medians run from 0.17 in nuclear, where more than half of the economies
share one value discussed below, and 0.20 in batteries and electric vehicles to 0.47 in oil.
Midstream and downstream, the pillar rests mainly on median-centred S-curves
(equation~\eqref{eq:scurve}), which place the median covered economy near 0.5 by construction;
the four midstream pairs scored on the Trade category alone (gas, coal, oil and nuclear) have
medians between 0.48 and 0.56. Two pairs depart from this pattern for reasons of construction.
Green Hydrogen midstream, scored on Foreign Dependency alone, gives 90 per cent of its 106
rankable economies the same value, 0.26, so that its box collapses to a line; Green Hydrogen
downstream, scored on Consumption alone, is bimodal, with an interquartile range from 0.10 to
0.72. The level of a pair is thus partly a property of its category bundle rather than of the
security of the economies in it, and comparisons of levels across pairs should be read with
Section~\ref{subsec:norm_pillar} in mind.

A variance decomposition over the 4{,}163 rankable cells makes the same point.
Technology-stage fixed effects explain 44 per cent of the variance of $\ES_{its}$, country fixed
effects alone 12 per cent, and the two together 60 per cent; country effects account for 28 per
cent of the variance that remains within pairs (the partial $R^2$ of country effects given pair
effects). Within stages, pair effects explain 33, 50 and 22 per cent of the upstream, midstream
and downstream variance, and country effects 51, 65 and 29 per cent of the within-pair
remainder. The pair effects are not an artefact of the cell-constant inputs: subtracting the
contributions of Energy Prices and of the world-concentration terms of Trade and Minerals Trade
from each cell leaves the pair-effect $R^2$ at 48 per cent. They reflect the different category
bundles, and in particular the weight of zero-rule registers in a pair. A cross-country ranking
of the pillar is therefore meaningful within a pair, and a comparison of one economy's levels
across pairs, which the Strategic Index makes, mixes economic content with the construction of
each pair.

\begin{figure}[tbp]
  \centering
  \includegraphics[width=\textwidth]{es_distribution}
  \caption{Distribution of the Energy Security pillar by technology and supply-chain stage,
  rankable cells. Each row summarises $\ES_{its}$ across the rankable economies of one
  technology-stage pair (between 68 and 195 economies; a cell is rankable if its economy is
  sovereign and at most half of its index weight is imputed). The box spans the interquartile
  range with the median marked, and the whiskers reach the most extreme value within 1.5
  interquartile ranges. Small gray dots are economies outside the coalition and coloured dots
  coalition members, in the stage colour (blue upstream, orange midstream, aqua downstream),
  both jittered vertically; the triangle marks the United States and the diamond China. Rows are
  ordered by their median within each stage. In Nuclear upstream, 137 of the 195 rankable
  economies have neither identified uranium resources nor mine production; the 105 of them
  whose trade in uranium ores and concentrates is not observed share one value, 0.172 (Energy
  Prices plus an imputed Trade category), and the other 32 lie between 0.145 and 0.201. In Green
  Hydrogen midstream 90 per cent of the economies share one value, 0.26, and the box collapses
  to a line. Geothermal upstream and midstream and Electric Grid downstream have no Energy
  Security score. The score lies in $[0,1]$, and higher values indicate a more secure position.
  Source: \texttt{energy\_security\_index} in \texttt{index\_outputs.rds}, final run;
  \texttt{scripts/paper/fig\_es.R}.}
  \label{fig:es_distribution}
\end{figure}

\paragraph{What drives the score.} Figure~\ref{fig:es_category_contrib} uses the exact category
decomposition of equation~\eqref{eq:cat_decomp} to show the mean share of each category in the
pillar, by stage, together with its mean effective weight. Upstream, Trade supplies 31 per cent
of the average score, Energy Prices 20 per cent, Minerals Trade 16 per cent and Foreign
Dependency 13 per cent, while Reserves and Production contribute 7.2 and 3.1 per cent against
mean effective weights of 20 and 19.5 per cent. The gap between share and weight is the zero
rule at work: for the typical economy the register categories are zero and add nothing, while
for the few holders and producers they are what places them at the top. Trade, by contrast,
gives almost every economy a positive score, and the upstream baskets extend it to four more
pairs, which makes it the largest upstream contributor. The upstream shares of Energy Prices and
Trade are raised further by Nuclear upstream, where the two together make the whole score of
every economy without uranium; without that pair they are 15 and 28 per cent. Midstream, Trade
supplies 73 per cent of the score and Foreign Dependency 26 per cent, that is, an economy's
share of world exports of the basket, its net import position, and its shares of refined mineral
supply and manufacturing capacity. Investment, defined midstream only for batteries, electric
vehicles, solar and wind, supplies 1.4 per cent against a mean effective weight of 7.0 per cent,
because the investment monitor records neither manufacturing investment nor capacity for most
economies (72 per cent of the rankable economies score zero in the median of these four pairs).
Downstream, Consumption (56 per cent) and Energy Access (24 per cent) dominate, with resource
potential, scored under Reserves, at 10 per cent and Investment at 9.8 per cent against a weight
of 19 per cent; the downstream pillar is mainly a measure of the depth and growth of domestic
deployment and use.

Part of every upstream score, and of every midstream score except those of Green Hydrogen, does
not vary across economies. Energy Prices is a world-level series, and the world-concentration
term of Trade takes one value per pair; the concentration term of Minerals Trade is constant for
all economies except two outside the World Bank list, Taiwan and R\'eunion. These cell-constant
inputs carry a mean effective weight of 0.21 upstream and 0.22 midstream, and none downstream,
but because the country-varying categories score low for most economies they supply on average
36 per cent of an upstream cell's score and 26 per cent of a midstream cell's score. They shift
the level of a pair without reordering the economies in it. Energy Prices illustrates the point:
its score ranges from 0.19 for gas to 0.81 for the electric grid, and it raises every economy's
upstream score by between 0.025 (batteries) and 0.088 (nuclear), leaving every within-pair rank
unchanged.

\begin{figure}[tbp]
  \centering
  \includegraphics[width=\textwidth]{es_category_contrib}
  \caption{Average composition of the Energy Security pillar by stage, rankable cells. For each
  cell the weighted category contributions $\omega_{kits}c_{kits}$ sum exactly to $\ES_{its}$
  (equation~\eqref{eq:cat_decomp}). Bars show the mean, across the rankable cells of a stage, of
  each category's share of the score, $\omega_{kits}c_{kits}/\ES_{its}$; the vertical tick shows
  the category's mean effective weight $\omega_{kits}$. A category absent from a cell counts as
  zero in both means, and a category with neither bar nor tick is not defined at that stage. A
  bar longer than its tick indicates a category whose score tends to exceed the cell's overall
  score; a bar much shorter than its tick, as for Reserves and Production upstream, indicates a
  category on which most economies score zero. The upstream Energy Prices and Trade shares
  include Nuclear upstream, where the two categories make the entire score of the 137 rankable
  economies without uranium (15 and 28 per cent without that pair). Labels give shares of at
  least 10 per cent. Sample: 1{,}124 upstream, 1{,}397 midstream and 1{,}637 downstream rankable
  cells (the five downstream cells with $\ES_{its}=0$ are dropped). Source:
  \texttt{energy\_security\_outputs\$category\_contributions} in \texttt{index\_outputs.rds};
  \texttt{scripts/paper/fig\_es.R}.}
  \label{fig:es_category_contrib}
\end{figure}

\paragraph{Uranium.} Nuclear upstream is scored on four categories: Energy Prices (the
volatility of the uranium price), Trade (the upstream basket of uranium ores and concentrates and
natural uranium compounds, HS 261210 and 284410, the latter moved from the midstream basket in
the final basket review), Reserves (identified recoverable uranium resources at a cost below
USD 260/kgU in the OECD-NEA/IAEA Red Book) and Production (uranium mine output from the same
source), with effective weights of one eighth, one quarter, one quarter and three eighths.
Figure~\ref{fig:es_nuclear_upstream} shows how the score is built. Of the 195 rankable
economies, 58 hold identified resources and 18 of these also produce. An economy with neither
scores the Energy Prices contribution, $0.704/8=0.088$, plus its Trade category. For the 105
economies whose uranium trade is not observed, Trade is imputed from the observed economies and
the score is one common value, 0.172; these economies, 54 per cent of the rankable set, share
rank 67. The other 32 economies without uranium have an observed Trade category and score
between 0.145 (Trinidad and Tobago) and 0.201 (Norway). The zeros on Reserves and Production and
an imputed Trade category give the economies without uranium imputed weight shares between one
quarter and one half, so all of them remain rankable. Because Production carries three eighths
of the weight and only producers score on it, the producers dominate the top of the ranking,
though no longer all of it: T\"urkiye, which holds resources but reports no mine output, ranks
18th (0.34), and the lowest-ranked producer is 26th. Kazakhstan leads (0.84), ahead of Canada
(0.74), Namibia (0.72) and Australia (0.72); China ranks 10th (0.56) and the United States 11th
(0.53), and 9 of the 15 highest-ranked economies are coalition members. Of the 57 members, 31
hold identified resources and 26 (46 per cent) do not, among them France, the United Kingdom,
South Korea, the Netherlands, Norway and Poland. The pair therefore separates uranium holders and
producers clearly but orders the economies without uranium only by their trade in uranium ores
and compounds, and most of them not at all. That matters for the Strategic Index: read as a
need, $1-\ES_{its}$, the common value becomes a uniform need of 0.83 for the 105 economies that
share it, a measurement of absence rather than of a gap that domestic supply could close. In the
Strategic Index sample, 86 of the 174 countries with a Nuclear upstream cell share that value,
and Section~\ref{subsec:si_allied} restricts allied upstream picks to economies with recorded
reserves or production of the stage's resources for this reason.

\begin{figure}[tbp]
  \centering
  \includegraphics[width=\textwidth]{es_nuclear_upstream}
  \caption{Composition of the Energy Security score in Nuclear upstream, rankable cells. Rows are
  the 15 highest-ranked of the 195 rankable economies, with their rank; coalition members are in
  bold. Each bar is $\ES_{its}$ split into its weighted category contributions
  (equation~\eqref{eq:cat_decomp}), which sum exactly to the score: Energy Prices (uranium price
  volatility, effective weight $1/8$, the same 0.088 for every economy), Trade (export share and
  net imports of uranium ores, concentrates and natural uranium compounds, HS 261210 and 284410,
  and the concentration of world exports; weight $1/4$), Reserves (identified recoverable
  resources below USD 260/kgU, OECD-NEA/IAEA Red Book; weight $1/4$) and Production (uranium mine
  output; weight $3/8$). The bottom row, below a hairline, stands for the 105 rankable economies
  with neither identified resources nor mine production whose uranium trade is not observed:
  their score is Energy Prices plus an imputed Trade category, 0.172. The other 32 rankable
  economies without uranium score between 0.145 and 0.201. Source:
  \texttt{energy\_security\_index} and
  \texttt{energy\_security\_outputs\$category\_contributions} in \texttt{index\_outputs.rds};
  \texttt{scripts/paper/fig\_es.R}.}
  \label{fig:es_nuclear_upstream}
\end{figure}

\paragraph{Correlation among categories.} Within pairs the categories are positively but far
from perfectly related. The median within-pair Spearman rank correlation between two categories
is 0.42 upstream, where only one of 78 pair-specific correlations is negative ($-0.01$), 0.33
midstream and 0.11 downstream. Upstream, the strongest link is between Reserves and Production,
with a median of 0.84 over the seven pairs that carry both; among the economies with a positive
score on both (a median of 16 per pair) it is still 0.87. The two categories, which carry mean
upstream effective weights of 20 and 19.5 per cent, describe nearly the same ordering of the
same small set of mining economies. Foreign Dependency correlates with Production at 0.76 and
with Reserves at 0.69. Midstream, Foreign Dependency and Trade correlate at a median of only
0.28, with a range across pairs from $-0.23$ to 0.48: the economies with large shares of
refined supply and manufacturing capacity are not always those with large export shares and
small net imports of the basket. Downstream the categories are close to unrelated: Energy Access
and Consumption correlate at 0.12, and resource potential correlates negatively with
Consumption ($-0.11$), so that the economies with the richest solar, wind and geothermal
resources are not those with the deepest installed base. Within a pair the ranking of the pillar
follows Energy Imports (median Spearman correlation with $\ES_{its}$ of 0.80), Minerals Trade
(0.79) and Reserves (0.79) upstream, Trade (0.94) midstream, and Investment (0.76) and
Consumption (0.74) downstream.

\paragraph{Leaders.} Table~\ref{tab:es_leaders} lists the three highest-scoring rankable
economies in each pair, the best-placed coalition member and the ranks of the United States and
China. China leads 11 of the 30 pairs: Batteries, Electric Vehicles and Wind upstream, where its
shares of world mineral supply, production and trade dominate the mineral categories; Batteries,
Electric Grid, Electric Vehicles, Solar and Wind midstream, where its shares of world
manufacturing capacity, refined supply and exports are largest \citep{iea2024etp,iea2025gcmo};
and Batteries, Nuclear and Wind downstream. No other economy leads more than two. Upstream, the
other leaders follow resource endowments: Saudi Arabia in oil, Qatar in gas, Australia in coal,
Kazakhstan in nuclear, South Africa in green hydrogen, whose mineral weights are dominated by the
platinum-group metals, and Chile in the electric grid and solar, in which copper carries most of
the demand weight among the minerals with reserve and production tables. In Electric Grid
upstream, whose trade basket holds copper and aluminium inputs rather than electricity, the
Democratic Republic of the Congo and Peru follow Chile. Midstream, the four pairs scored on
Trade alone are led by economies with large export shares of the basket and small net imports
relative to GDP: Poland in coal, South Korea in gas and oil, and Russia in nuclear, whose basket
now holds enriched uranium, reactor components and steam-cycle equipment rather than natural
uranium. Green Hydrogen midstream, in which 95 of the 106 rankable economies share one value, is
led by Indonesia. Downstream the leaders follow domestic deployment and resource potential: Qatar
in oil, Kuwait in gas, Indonesia in coal, Saudi Arabia in solar power, the United States in
geothermal power and green hydrogen, and T\"urkiye in electric vehicles, a pair that rests on the
Investment category alone and ranks 121 economies. A coalition member is first in 16 of the 30
pairs and in the top three of all 30.

Figure~\ref{fig:es_top} shows the ten highest-ranked rankable economies in three illustrative
pairs. In Oil upstream Saudi Arabia leads (0.81) among the large net exporters with large proved
reserves; two coalition members, Saudi Arabia and Canada, are in the top ten, the United States
is 12th (0.71) and China, a large net importer, 34th (0.55). Every economy of the pair receives
0.16 from the cell-constant inputs, which therefore account for none of the gap between the mean
score of the top ten and the mean over all 97 rankable economies; Reserves accounts for 58 per
cent of that gap, Energy Imports for 28 per cent and the country-varying part of Trade for 13
per cent. China's position is that of an importer: its Reserves and Trade scores (0.64 and 0.79)
are not far from those of the United States (0.69 and 0.87), but its Energy Imports score is
0.25 against 0.68. In the two midstream pairs China leads by a wide margin, with 0.77 in
batteries against 0.56 for the second-placed United States, and 0.82 in solar against 0.67 for
India, and coalition members fill all nine remaining top-ten places of each pair, among them the
Southeast Asian manufacturing economies that the supplier and agreement routes bring into the
coalition (Indonesia and Malaysia in batteries; Thailand, Malaysia, Viet Nam and Indonesia in
solar). Investment accounts for 62 per cent of the gap between the top ten and all rankable
economies in batteries and 58 per cent in solar. It is a zero-rule category on which the mean
economy scores 0.13 and 0.09 and the top ten 0.61 and 0.58: the economies near the top are those
that the investment monitor records as investing in, operating or building manufacturing
capacity. China's lead over the second-placed economy comes from Foreign Dependency, that is,
from its shares of refined mineral supply and of manufacturing capacity (a weighted contribution
of 0.37 in batteries against 0.21 for the United States, and 0.40 in solar against 0.28 for
India); on Investment the United States scores higher than China in batteries (0.86 against
0.80).

% Generated by scripts/paper/fig_es.R; do not edit by hand.
\begin{table}[tbp]
\centering
\begin{threeparttable}
\caption{Leading economies in each technology-stage pair of the Energy Security pillar (rankable cells)}
\label{tab:es_leaders}
\footnotesize
\begin{tabular}{@{}lrlllcrr@{}}
\toprule
Technology & $n$ & First & Second & Third & Best member (rank) & USA & CHN \\
\midrule
\multicolumn{8}{@{}l}{\textit{Upstream}}\\
Batteries & 106 & CHN 0.70 & AUS 0.60 & RUS 0.56 & AUS (2) & 9 & 1 \\
Coal & 96 & AUS 0.72 & USA 0.72 & RUS 0.71 & AUS (1) & 2 & 8 \\
Electric Grid & 106 & CHL 0.88 & COD 0.74 & PER 0.72 & CHL (1) & 5 & 4 \\
Electric Vehicles & 107 & CHN 0.70 & AUS 0.61 & RUS 0.58 & AUS (2) & 5 & 1 \\
Gas & 98 & QAT 0.75 & RUS 0.73 & USA 0.69 & QAT (1) & 3 & 33 \\
Green Hydrogen & 106 & ZAF 0.66 & RUS 0.64 & IDN 0.63 & ZAF (1) & 7 & 8 \\
Nuclear & 195 & KAZ 0.84 & CAN 0.74 & NAM 0.72 & CAN (2) & 11 & 10 \\
Oil & 97 & SAU 0.81 & VEN 0.79 & IRQ 0.78 & SAU (1) & 12 & 34 \\
Solar & 106 & CHL 0.75 & CHN 0.65 & PER 0.64 & CHL (1) & 5 & 2 \\
Wind & 107 & CHN 0.73 & AUS 0.69 & USA 0.63 & AUS (2) & 3 & 1 \\
\addlinespace
\multicolumn{8}{@{}l}{\textit{Midstream}}\\
Batteries & 106 & CHN 0.77 & USA 0.56 & CAN 0.53 & USA (2) & 2 & 1 \\
Coal & 190 & POL 0.80 & MOZ 0.80 & CHN 0.79 & POL (1) & 10 & 3 \\
Electric Grid & 106 & CHN 0.78 & JPN 0.70 & CHL 0.64 & JPN (2) & 22 & 1 \\
Electric Vehicles & 105 & CHN 0.80 & ESP 0.67 & HUN 0.66 & ESP (2) & 4 & 1 \\
Gas & 192 & KOR 0.84 & CHN 0.83 & ITA 0.81 & KOR (1) & 8 & 2 \\
Green Hydrogen & 106 & IDN 0.61 & CHN 0.60 & RUS 0.49 & IDN (1) & 12$^{95}$ & 2 \\
Nuclear & 188 & RUS 0.78 & FRA 0.74 & NLD 0.73 & FRA (2) & 30 & 5 \\
Oil & 192 & KOR 0.78 & CHN 0.77 & GAB 0.74 & KOR (1) & 11 & 2 \\
Solar & 106 & CHN 0.82 & IND 0.67 & USA 0.59 & IND (2) & 3 & 1 \\
Wind & 106 & CHN 0.80 & POL 0.59 & USA 0.56 & POL (2) & 3 & 1 \\
\addlinespace
\multicolumn{8}{@{}l}{\textit{Downstream}}\\
Batteries & 186 & CHN 0.67 & DEU 0.66 & USA 0.65 & DEU (2) & 3 & 1 \\
Coal & 186 & IDN 0.74 & CHN 0.74 & JPN 0.69 & IDN (1) & 9 & 2 \\
Electric Vehicles & 121 & TUR 0.90 & CHN 0.88 & VNM 0.82 & TUR (1) & 4 & 2 \\
Gas & 186 & KWT 0.71 & QAT 0.70 & KOR 0.67 & QAT (2) & 6 & 24 \\
Geothermal & 146 & USA 0.83 & IDN 0.64 & JPN 0.58 & USA (1) & 1 & 9 \\
Green Hydrogen & 186 & USA 0.74 & KOR 0.74 & ALB 0.72$^{49}$ & USA (1) & 1 & 63 \\
Nuclear & 68 & CHN 0.72 & FRA 0.71 & JPN 0.66 & FRA (2) & 17 & 1 \\
Oil & 186 & QAT 0.91 & OMN 0.85 & ARE 0.83 & QAT (1) & 23 & 186 \\
Solar & 191 & SAU 0.70 & AUS 0.70 & CHN 0.69 & SAU (1) & 4 & 3 \\
Wind & 186 & CHN 0.67 & CAN 0.66 & USA 0.65 & CAN (2) & 3 & 1 \\
\bottomrule
\end{tabular}
\begin{tablenotes}[flushleft]\footnotesize
\item \textit{Notes:} $\ES_{its}$ for the eleven headline technologies, final run. Only rankable cells enter: a sovereign economy with at most half of the cell's index weight imputed. Economies are identified by ISO3 code and followed by their score; $n$ is the number of rankable economies in the pair. ``Best member'' is the highest-ranked of the 57 coalition members, with its rank in parentheses; the USA and CHN columns give the ranks of the United States and China. A rank is one plus the number of rankable economies with a strictly higher score, so tied economies share a rank, and a superscript gives the number of economies that share it (the economy listed is the first of them alphabetically). Geothermal upstream, Geothermal midstream and Electric Grid downstream have no Energy Security score. Source: \texttt{energy\_security\_index} in \texttt{index\_outputs.rds}; \texttt{scripts/paper/fig\_es.R}.
\end{tablenotes}
\end{threeparttable}
\end{table}


\begin{figure}[tbp]
  \centering
  \includegraphics[width=\textwidth]{es_top}
  \caption{The ten highest-ranked economies in three illustrative technology-stage pairs, Oil
  upstream, Batteries midstream and Solar midstream, rankable cells. Each dot is the economy's
  $\ES_{its}$ (scale 0 to 1), and the number before its name is its rank among the rankable
  economies of the pair, one plus the number of economies with a strictly higher score.
  Coalition members are in bold and drawn in the stage colour; other economies are gray; the
  triangle marks the United States and the diamond China, which are appended below a hairline
  with their rank when they fall outside the top ten. The vertical line is the median over the
  rankable economies of the pair (97, 106 and 106 economies). The shaded band from zero is the
  cell-constant part of the score, the same for every economy of the pair: Energy Prices and the
  world-concentration term of Trade in Oil upstream (0.16), the latter alone in the two midstream
  pairs (0.03 and 0.06). Ties are listed alphabetically. Source:
  \texttt{energy\_security\_index} and
  \texttt{energy\_security\_outputs\$variable\_contributions} in \texttt{index\_outputs.rds};
  \texttt{scripts/paper/fig\_es.R}.}
  \label{fig:es_top}
\end{figure}

\paragraph{The coalition.} The 57 coalition members of Section~\ref{sec:coalition} score above
the other economies at every stage: their mean $\ES_{its}$ over the rankable cells is 0.34
upstream, 0.49 midstream and 0.48 downstream, against 0.29, 0.44 and 0.42 for non-members, and
their mean within-pair percentile rank is 0.59, 0.72 and 0.62, against 0.47, 0.41 and 0.50. The
two tiers differ by stage. The partners, several of which enter the coalition as major
suppliers, have the higher upstream mean percentile (0.63 against 0.56 for the core), while the
core is ahead midstream (0.78 against 0.64) and downstream (0.64 against 0.59). The 21 EU members
of the coalition are weaker upstream than the other members (mean percentile 0.50 against 0.64)
and stronger midstream (0.79 against 0.67). The coalition's advantage is unsurprising, since it
consists mainly of large or resource-rich economies, and it says nothing about whether members'
strengths complement one another, which Section~\ref{sec:network} takes up.

Figure~\ref{fig:es_heatmap} shows the members' scores across all pairs. Ordered by mean
within-pair percentile over their rankable pairs, the United States (0.95), Canada (0.89),
Australia (0.82), India (0.79) and Poland (0.77) lead, and Qatar (0.44), Ghana (0.45), Gabon
(0.46), Namibia (0.46) and Luxembourg (0.47) trail. The upstream columns separate the resource
holders, with the United States, Australia, Canada, Indonesia and the Democratic Republic of the
Congo the members with the highest upstream mean percentiles, from the industrial economies.
Midstream the ordering changes: Japan has the highest midstream mean percentile of any member
(0.95), ahead of Germany, the United States, India and South Korea, which reflects manufacturing
and export positions rather than endowments. Downstream the ordering is less systematic, as the
low partial $R^2$ of country effects downstream and the weak correlation among the downstream
categories would lead one to expect. Of the 1{,}710 member cells in the 30 scored pairs,
1{,}606 are rankable, 96 are scored but not rankable and 8 are not scored. The non-rankable
cells concentrate in five partners, Namibia, Costa Rica, the Dominican Republic, Guatemala and
Ghana, which are rankable in only 13 to 15 of their 28 or 29 scored pairs, and in Canada (23 of
29), so that their rows rest on fewer comparisons than the others.

\begin{figure}[tbp]
  \centering
  \includegraphics[width=\textwidth]{es_heatmap}
  \caption{Energy Security pillar of the 57 coalition members across the 33 technology-stage
  sectors. Rows are grouped into the ten core members outside the European Union, the 21 core
  members in the EU, which the network of Section~\ref{sec:network} treats as one node together
  with six smaller EU economies that are not members in their own right, and the 26 partners;
  within each group, rows are ordered by the member's mean within-pair percentile rank over its
  rankable pairs (highest at the top). China is shown below for reference. Columns are grouped
  by stage. Colour gives $\ES_{its}$ on a sequential scale from 0 to 1, darker meaning more
  secure, for rankable cells; gray cells are scored but not rankable (a sovereign economy with
  more than half of the cell's index weight imputed); a cross marks a member with no score in a
  pair; the three sectors without an Energy Security score (Geothermal upstream and midstream,
  Electric Grid downstream) are left blank. Source: \texttt{energy\_security\_index} and
  \texttt{coalition\_members} in \texttt{index\_outputs.rds}; \texttt{scripts/paper/fig\_es.R}.}
  \label{fig:es_heatmap}
\end{figure}

\paragraph{The United States, China and Taiwan.} The two largest economies rank near the top of
most pairs, but for different reasons. The United States has the highest mean percentile rank
of any coalition member, 0.95, with a place in the top ten in 22 of the 30 pairs; it leads two,
Geothermal and Green Hydrogen downstream. Its upstream position rests on fossil and mineral
resources (second in coal, third in gas and wind upstream), and its weakest positions are 30th
in Nuclear midstream, 23rd in Oil downstream, 22nd in Electric Grid midstream and 17th in
Nuclear downstream. China's mean percentile rank is 0.92, but it is in the top ten in 25 pairs
and first in 11, and its midstream mean percentile is 0.99, with a place in the top ten of
every midstream pair. Its weak positions are those of a large fossil-fuel importer, 33rd in gas
and 34th in oil upstream and 24th in gas downstream, and two downstream pairs in which the
deployment data place it low: 63rd in green hydrogen and last of 186 in oil. In the oil pair
Consumption carries five-ninths of the weight, no rankable economy has a lower projected growth
of oil-fired capacity, and only one a lower current level; the measure rewards a large and
expanding domestic oil-fired fleet. The United States ranks above China in 8 of the 30 pairs,
all of them upstream or downstream fossil-fuel and resource pairs (coal, gas, green hydrogen and
oil upstream; gas, geothermal, green hydrogen and oil downstream) and none of them midstream.
The contrast is only partly one between an endowment-based and a manufacturing-based position,
since China also leads three upstream mineral pairs; it mirrors, rather than explains, the
change in the pillar's leaders from stage to stage that follows from the different category
bundles of the stages. Taiwan, which the trade themes now carry, is scored in 28 cells and
rankable in 19. Its profile is that of a manufacturing economy without resources: a mean
within-pair percentile of 0.75 across the ten midstream pairs, with ranks between 27th and 61st
on the Trade category, against 0.01 across its seven rankable upstream pairs, near the bottom of
each, and no place in the top ten of any pair.

\paragraph{Caveats.} Five features of the construction qualify these results. First,
technology-stage effects dominate the variance of the pillar, so it is a cross-country measure
within pairs and only partly comparable across them. Second, the cell-constant inputs supply on
average more than a third of an upstream cell's score and about a quarter of a midstream cell's
score without informing any ranking. Third, several pairs rank few distinct values. In Green
Hydrogen midstream the rank of 12 that the United States holds is shared with 94 other
economies, and in Green Hydrogen downstream 49 economies share third place at 0.72. In the
Strategic Index sample five sectors have at least a quarter of their countries at one Energy
Security value: Green Hydrogen midstream (34 of 40 countries), Nuclear upstream (86 of 174),
Green Hydrogen downstream (45 of 158), Nuclear downstream (19 of 67) and Batteries downstream
(38 of 138). Under the need reading of the Strategic Index such a mass point turns into a
uniformly high need for every country at it (Section~\ref{sec:strategic}). Fourth, the register
categories count as zero any holder or producer that a source does not quantify: holders that
the USGS Mineral Commodity Summaries list without a figure (not available, withheld or small)
score zero in the technology roll-ups of Reserves. Fifth, the downstream pillar treats a deeper
and faster-growing domestic use of a carrier as security; a reader who regards high consumption
of a fuel as exposure would reverse the downstream reading for the fossil fuels.
